% !TEX program = lualatex
% CV français ciblé pour Inria Startup Studio / Percolia.
% Compilation : lualatex -interaction=nonstopmode -halt-on-error CV_Alban_Hauseux_Inria_Startup_Studio_2026.tex
\documentclass[10pt,a4paper]{article}

\usepackage{fontspec}
\setsansfont{TeX Gyre Heros}[Ligatures=TeX]
\renewcommand{\familydefault}{\sfdefault}
\usepackage[left=1.35cm,right=1.35cm,top=0.85cm,bottom=0.85cm]{geometry}
\usepackage{microtype}
\usepackage{xcolor}
\usepackage{titlesec}
\usepackage{enumitem}
\usepackage{tabularx}
\usepackage{hyperref}

\definecolor{darkblue}{RGB}{20,52,80}
\definecolor{midblue}{RGB}{53,94,128}
\definecolor{softgray}{RGB}{78,78,78}

\hypersetup{
  unicode=true,
  colorlinks=true,
  urlcolor=darkblue,
  linkcolor=darkblue,
  pdfauthor={Alban Hauseux},
  pdftitle={Curriculum vitae -- Alban Hauseux -- Inria Startup Studio},
  pdfsubject={Candidature Inria Startup Studio -- projet Percolia},
  pdfkeywords={Inria Startup Studio, Percolia, porteur business, produit, Python, SQL, finance quantitative}
}
\urlstyle{same}

\titleformat{\section}
  {\normalsize\bfseries\color{darkblue}}
  {}{0pt}{\MakeUppercase}[\vspace{-0.32em}\color{midblue}\rule{\textwidth}{0.45pt}]
\titlespacing*{\section}{0pt}{0.52em}{0.24em}

\setlength{\parindent}{0pt}
\setlength{\tabcolsep}{0pt}
\setlist[itemize]{
  leftmargin=2.68cm,
  label=\textcolor{midblue}{\raisebox{0.12ex}{\scriptsize$\bullet$}},
  labelsep=0.42em,
  itemsep=0.03em,
  topsep=0.05em,
  parsep=0pt,
  partopsep=0pt
}

\newcommand{\cventry}[4]{%
  \begin{tabularx}{\textwidth}{@{}p{2.32cm}@{\hspace{0.22cm}}X@{}}
    \textcolor{softgray}{#1} & \textbf{#2}\if\relax\detokenize{#3}\relax\else\hfill\textcolor{softgray}{#3}\fi\\[-0.15ex]
    & \textbf{#4}
  \end{tabularx}\vspace{-0.25em}
}

\newcommand{\compactentry}[2]{%
  \begin{tabularx}{\textwidth}{@{}p{2.32cm}@{\hspace{0.22cm}}X@{}}
    \textcolor{softgray}{#1} & #2
  \end{tabularx}\par
}

\begin{document}
\pagestyle{empty}
\frenchspacing

% ==================== EN-TÊTE ====================
\begin{center}
  {\LARGE\bfseries Alban \textsc{Hauseux}}\\[0.22em]
  {\large\bfseries\color{darkblue}Porteur business de Percolia}\\[0.12em]
  {\normalsize Développement commercial, audits clients et intégration logicielle}\\[0.32em]
  Paris
  \enspace\textcolor{midblue}{\textbullet}\enspace
  +33~(0)6~87~59~29~26
  \enspace\textcolor{midblue}{\textbullet}\enspace
  \href{mailto:alban.hau@gmail.com}{alban.hau@gmail.com}
\end{center}
\vspace{-0.55em}

% ==================== PROFIL ====================
\section*{Profil}
Ingénieur télécom et diplômé du master Économie et Finance de l'Université Paris-Dauphine
PSL, expérimenté en développement d'outils quantitatifs, vente et exécution de produits
financiers, contrôle de modèles et coordination Front Office--IT. Dans \textbf{Percolia},
mon rôle prévu est de structurer le développement commercial et les audits clients, puis
d'accompagner l'intégration logicielle aux côtés de Louis Hauseux, porteur de la technologie.
Candidature pour un engagement à temps plein au sein d'Inria Startup Studio.

% ==================== PERCOLIA ====================
\section*{Projet entrepreneurial}
\cventry{2026}{Percolia -- projet candidat à Inria Startup Studio}{}{Porteur business -- responsabilités prévues}
\begin{itemize}
  \item Structurer la découverte client, le démarchage et les audits comparatifs payants sur
  le cas d'usage prioritaire des nuages de points 3D LiDAR, puis convertir les cas validés en
  API ou licences logicielles récurrentes.
  \item Piloter l'interface client--R\&D : données et besoins métier, spécifications, tests
  d'acceptation, restitution des résultats et intégration ; contribuer au développement en
  Python et SQL.
\end{itemize}

% ==================== EXPÉRIENCE ====================
\section*{Expériences professionnelles}
\cventry{2023--auj.}{HSBC}{Paris}{Assistant Trader Repo (depuis 2025) ; Rates Products (2023--2024)}
\begin{itemize}
  \item Développer des outils d'analyse et d'aide à la décision en Python, SQL et VBA ;
  produire notamment le calcul du P\&L de la salle des marchés.
  \item Assurer le support du desk et l'interface Front Office--Middle/Back Office--IT ;
  gérer le booking de swaps, FRA, forward bonds et swaps d'inflation zéro-coupon.
  \item Reconcevoir un processus d'allocation du P\&L sur un périmètre de taux souverains
  européens : spécifications, procédures et tests menés avec l'IT.
\end{itemize}

\cventry{2021--2022}{Société Générale}{Luxembourg}{Cross-Asset Structured Products Sales and Execution Trader}
\begin{itemize}
  \item Exécuter les ordres sur les marchés gris et secondaire, contrôler les prix des
  émetteurs et proposer des produits structurés adaptés aux besoins des clients.
  \item Développer en Python un outil de contrôle des prix connecté à Bloomberg et au système
  de booking.
\end{itemize}

\cventry{2020--2021}{HSBC}{Paris}{Assistant Trader Exotic Equity Derivatives -- développement Python \& VBA (alternance)}
\begin{itemize}
  \item Développer et améliorer les outils du desk : financement d'autocallables, couverture
  du vega, suivi du P\&L et backtests de stratégies.
\end{itemize}

\cventry{2019--2020}{Crédit Mutuel}{Paris}{Audit quantitatif -- Inspection générale groupe}
\begin{itemize}
  \item Contribuer à des missions sur les modèles de crédit, les critères de défaut bâlois
  et la réconciliation de bases de données opérationnelles et décisionnelles.
\end{itemize}

\cventry{2018--2019}{Banque de France--ACPR}{Paris}{Contrôle des modèles internes Solvabilité~II}
\begin{itemize}
  \item Contribuer aux contrôles sur place d'assureurs : risques de marché et d'assurance,
  conformité prudentielle, modèles de taux et modèles actions.
  \item Implémenter en Python un générateur de scénarios économiques mobilisant modèles de
  taux et actions, équations différentielles stochastiques, optimisation et arbres multinomiaux.
\end{itemize}

% ==================== COMPÉTENCES ====================
\section*{Compétences}
\compactentry{Logiciel \& data}{\textbf{Python, SQL, VBA} ; API Bloomberg ; automatisation,
analyse de données, backtests, scénarios économiques et outils de P\&L.}
\compactentry{Produit \& projet}{Spécifications fonctionnelles, procédures, plans de test,
coordination métiers--IT, restitution d'audits et intégration d'outils décisionnels.}
\compactentry{Quantitatif}{Repo, taux, dérivés actions, produits structurés ;
exécution, relation client, pricing, modèles de risque et réglementation prudentielle.}

% ==================== FORMATION ====================
\section*{Formation}
\compactentry{2020--2021}{\textbf{Master Économie et Finance}, Université Paris-Dauphine PSL -- Paris.}
\compactentry{2015--2018}{\textbf{Diplôme d'ingénieur}, IMT Atlantique (ex-Télécom Bretagne) -- Brest.}
\compactentry{2012--2015}{Classes préparatoires scientifiques, lycée Saint-Louis -- Paris.}

\end{document}
